% 04a2_release_gap.tex -- the release transient against the survivors' largest gap (twin, 1 m and 0.55 m).
\subsection{What sets the release transient}
\label{sec:res_release_gap}

The rig's slowest release left the largest gap between survivors (Section~\ref{sec:res_reconfig}). To isolate that
variable, the twin flies the same unannounced release from four carriers in hover, with the carriers placed so that the
three survivors leave a largest gap of 120 to 180$^\circ$ around the ring. Everything else is held fixed. We report the
rise of the ring's tilt from the release to 15\,s after it, cut at the moment the ring is dropped or a fault stops the
flight. One arm at 150$^\circ$ feeds the survivors' trackers the planned cable pull instead of the measured one
(Section~\ref{sec:planner}). The 0.55\,m column is one run a cell.

Table~\ref{tab:release_gap} gives the results. The tilt rise is steep in the gap: from 120 to 150$^\circ$ it more than
doubles, at 165$^\circ$ it crosses the 30$^\circ$ line at which the ring is dropped and every carrier lands (2 of 2),
and at 180$^\circ$, with the survivors on one half of the rim, it passes 60$^\circ$ and the fleet is stopped by the
tilt fault. Every neighbouring pair of gaps is separated by more than the spread of its repeats at 1\,m.
The longer cable changes the boundary only slightly: at 0.55\,m the 165$^\circ$ gap rose just under the drop line
and the ring was held (one run). At 150$^\circ$ the measured pull holds the rise to about half of the rise with the planned pull at 0.55\,m and
to about two thirds at 1\,m, where the planned arm stays just under the drop line instead of crossing it.

\begin{table}[t]
\centering
\footnotesize
\setlength{\tabcolsep}{3pt}
\caption{Rise of the ring's tilt after an unannounced release from four carriers in hover (twin), against the
survivors' largest gap. Repeats as a\,/\,b; $\geq$: cut at the drop or the fault. Measured cable pull unless marked.}
\label{tab:release_gap}
\begin{tabular}{lccl}
\toprule
Largest gap & Tilt rise, 1\,m & Tilt rise, 0.55\,m & Outcome, 1\,m \\
($^\circ$) & ($^\circ$) & ($^\circ$) & \\
\midrule
120 & 6.90\,/\,7.45\,/\,7.50 & 6.40 & held \\
135 & 10.42\,/\,11.02 & 8.44 & held \\
150 & 18.17\,/\,17.58 & 16.82 & held \\
150, planned pull & 29.62\,/\,27.10 & $\geq$29.54 & held (0.55\,m: dropped) \\
165 & $\geq$30.49\,/\,$\geq$29.84 & 28.02 & dropped (0.55\,m: held) \\
180 & $\geq$62.46\,/\,$\geq$61.91 & $\geq$61.22 & tilt fault \\
\bottomrule
\end{tabular}
\end{table}

The controller cannot undo a bad geometry: once the survivors leave a wide gap, the ring has lost its support on that
side, and no redistribution of pull among them restores it within the drop line. The leaving order is therefore part
of the design. An order that keeps the survivors' gaps small keeps the transient small; opposite pairs are the simplest
such order. The $8\to4\to8$ chain of Fig.~\ref{fig:s_chain} releases its carriers in opposite pairs, so each release
leaves gaps of at most 90$^\circ$.
